% SPDX-License-Identifier: Apache-2.0
\section{Values}
\label{sec:e-values}

Rust's integers are the wrong types for hardware, and the reason is not
width. It is that \code{u32} has no way to say that nobody has driven this
wire yet.

\code{crate::types} therefore defines its own, paralleling the primitives
rather than reusing them.

\subsection{Two values and nine}

\code{Bit} is two valued, and it is what synthesis maps. Its default is
\code{Zero}, because a synthesised register leaves reset at a defined
value.

\code{Logic} is nine valued, in the order IEEE 1164 gives
\code{std\_ulogic}. Its default is \code{U}, uninitialised, and that
difference is the whole reason the type exists.

\begin{lstlisting}[caption={The two value types. An undriven wire is \code{U}, not zero.},label={lst:bitlogic}]
pub enum Bit { Zero, One }

pub enum Logic {
    U,        // uninitialised
    X,        // forcing unknown
    Zero, One,
    Z,        // high impedance
    W,        // weak unknown
    L, H,     // weak 0, weak 1
    DontCare,
}
\end{lstlisting}

Narrowing is fallible, and that is the point. \code{Logic::to\_bit}
returns nothing for \code{U}, \code{X} or \code{Z}, so a design that turns
an unknown into a number has to say so rather than do it quietly.

\subsection{Resolution, not a rule of our own}

Two drivers on one wire are resolved by the IEEE 1164 table, which the
library implements rather than paraphrases.

\begin{lstlisting}[caption={Resolution. Contention gives \code{X}; a pull-up against high impedance gives \code{H}.},label={lst:resolve}]
Logic::One.resolve(Logic::Zero)   // X
Logic::Z.resolve(Logic::H)        // H
\end{lstlisting}

Reusing the standard table costs nothing and means an engineer who knows
VHDL already knows what this does.

\subsection{Numbers}

\code{U<N>} and \code{I<N>} are the numeric vectors, and
\code{logic::Vec<N>} is a vector of \code{Logic} for when unknowns must
propagate. It lives in its own module so that the type is just
\code{Vec}: \code{logic::Vec<8>} says what it is without repeating the
word, and \code{std::vec::Vec} is untouched, because nothing is imported
unqualified. Converting between them is where the two worlds meet:
\code{logic::Vec::to\_u} returns nothing unless every bit is defined.

\begin{lstlisting}[caption={An undriven vector is unknown, and says so.},label={lst:undriven}]
let v = logic::Vec::<8>::default();   // eight U's
assert!(v.to_u().is_none());
\end{lstlisting}

Same-width arithmetic, resizing to a stated width and slicing all compile
on stable, because each result width is a standalone const parameter. So
does a multiply whose result width is \emph{stated}, \code{a.mul::<64>(b)},
because \code{64} is standalone too. What does not compile is a multiply
whose result width is \emph{computed}, \code{U<\{A + B\}>}, for the
reason Section~\ref{sec:e-probes} gives: that is the one thing stable
Rust will not accept, so the width is written by the caller rather than
derived.

A number is written as a number. \code{U<N>} converts from every unsigned
primitive and from \code{i32}, and its operators and every drive take
anything that converts, so \code{acc + 1}, \code{n == 0} and
\code{reg.set(0)} say what they mean. The operators are Rust's,
\code{+}, \code{-}, \code{\&}, \code{|}, \code{\^{}}, \code{!},
\code{<<}, \code{>>} and the compares, through the \code{std::ops}
traits; a compare yields a \code{bool}, since \code{PartialEq} allows
nothing else, and a \code{bool} converts to a \code{Bit} and back, so
a condition is either. \code{i32} is there because an
integer literal with no other constraint is an \code{i32}, and a literal
is the common case; a negative one is a bug, and debug builds say so.
